\documentclass[a4paper,11pt]{article}

\usepackage{amsmath,amssymb,amsfonts}
\usepackage{fontspec}
\usepackage{libertinus-otf}
\usepackage[french]{babel}
\usepackage{array}
\usepackage{booktabs}
\usepackage{longtable}
\usepackage{geometry}
\usepackage{multirow}
\usepackage[table]{xcolor}
\usepackage{colortbl}
\usepackage{fancyhdr}
\usepackage{lastpage}
\usepackage{enumitem}
\usepackage{eurosym}
\usepackage{tikz}
\usepackage[most]{tcolorbox}
\usepackage{tabularx}

\geometry{
    a4paper,
    left=1.1cm,
    right=1.1cm,
    top=1.6cm,
    bottom=1.5cm,
    headheight=20pt,
    headsep=10pt,
    footskip=20pt
}

\setlength{\parindent}{0pt}
\setlength{\tabcolsep}{4.5pt}
\renewcommand{\arraystretch}{1.18}
\setlength{\aboverulesep}{1.5pt}
\setlength{\belowrulesep}{1.5pt}

\definecolor{bleuprimaire}{RGB}{20, 52, 98}
\definecolor{bleusecondaire}{RGB}{31, 95, 168}
\definecolor{bleufond}{RGB}{238, 244, 251}
\definecolor{grisaccent}{RGB}{100, 116, 139}
\definecolor{vertvrai}{RGB}{16, 120, 60}
\definecolor{rougefaux}{RGB}{195, 20, 20}
\definecolor{filetgris}{RGB}{220, 226, 235}

\tikzset{
    tabcadre/.style={draw=bleuprimaire, line width=0.8pt},
    tabligne/.style={draw=bleuprimaire!45, line width=0.6pt},
    tabrep/.style={
        draw=rougefaux!90,
        fill=rougefaux!8,
        text=rougefaux,
        font=\bfseries\footnotesize,
        rounded corners=2.5pt,
        inner sep=1.8pt,
        line width=0.6pt
    },
    tabfleche/.style={
        ->,
        >=stealth,
        thick,
        draw=bleusecondaire,
        shorten <= 3pt,
        shorten >= 3pt
    }
}

\newcommand{\pts}[1]{%
    \tcbox[on line, size=small, colback=bleufond, colframe=bleusecondaire!50,
           arc=2.5mm, boxrule=0.6pt, top=1.2pt, bottom=1.2pt, left=4pt, right=4pt]{%
        \small\bfseries\textcolor{bleuprimaire}{#1}%
    }%
}

\newcommand{\subpt}[1]{%
    \tcbox[on line, size=small, colback=bleufond, colframe=bleusecondaire!30,
           arc=2mm, boxrule=0.5pt, top=1pt, bottom=1pt, left=3pt, right=3pt]{%
        \tiny\bfseries\textcolor{bleuprimaire}{#1}%
    }%
}

\newcommand{\repbox}[1]{%
    \tcbox[on line, size=small, colback=rougefaux!6, colframe=rougefaux!90,
           arc=1.2mm, boxrule=0.7pt, top=1pt, bottom=1pt, left=3.5pt, right=3.5pt]{%
        \bfseries\textcolor{rougefaux}{#1}%
    }%
}

\newcommand{\vect}[1]{\overrightarrow{#1}}

\pagestyle{fancy}
\fancyhf{}
\lhead{\small\bfseries\textcolor{bleuprimaire}{BAC BLANC} \textcolor{grisaccent}{-- Terminale Spécialité Mathématiques}}
\rhead{\small\bfseries\textcolor{bleusecondaire}{Barème et Corrigé détaillé}}
\cfoot{\small\textcolor{grisaccent}{Page \thepage\ sur \pageref{LastPage}}}

\renewcommand{\headrulewidth}{0.5pt}
\renewcommand{\headrule}{\hbox to\headwidth{\color{bleusecondaire!60}\leaders\hrule height \headrulewidth\hfill}}
\renewcommand{\footrulewidth}{0.3pt}
\renewcommand{\footrule}{\hbox to\headwidth{\color{filetgris}\leaders\hrule height \footrulewidth\hfill}}

\begin{document}

\begin{tcolorbox}[
    enhanced,
    colback=bleufond,
    colframe=bleuprimaire,
    arc=1.6mm,
    boxrule=0.8pt,
    top=4pt, bottom=4pt, left=8pt, right=8pt
]
    \textbf{\Large\textcolor{bleuprimaire}{Grille d'évaluation chiffrée et Corrigé détaillé}} \hfill \textbf{\Large\textcolor{bleuprimaire}{Total : 20 points}} \\[0.15em]
    \textbf{\textcolor{bleusecondaire}{Session Mars 2027}} \hspace{0.2cm}\textcolor{grisaccent}{\textbullet}\hspace{0.2cm} \textcolor{grisaccent}{Bac Blanc Mathématiques} \hfill \small\textcolor{grisaccent}{Durée : 4 heures \textbullet\ Calculatrice autorisée}
\end{tcolorbox}

\vspace{0.10cm}

\begin{center}
\renewcommand{\arraystretch}{1.4}
\begin{tabularx}{\textwidth}{|X|c|c|c|c|c|}
    \hline
    \textbf{Exercice} & \textbf{Exercice 1} & \textbf{Exercice 2} & \textbf{Exercice 3} & \textbf{Exercice 4} & \textbf{Total} \\ \hline
    \textbf{Thème} & Analyse et Suites & Géométrie espace & Probabilités & Fonctions & --- \\ \hline
    \textbf{Barème} & 4 points & 5 points & 6 points & 5 points & \textbf{20 points} \\ \hline
\end{tabularx}
\end{center}

\vspace{0.2cm}

\begin{longtable}{
    >{\centering\arraybackslash\bfseries}p{1.1cm}
    >{\raggedright\arraybackslash\bfseries}p{4.1cm}
    !{\color{filetgris}\vrule width 0.7pt}
    >{\raggedright\arraybackslash}p{11.1cm}
    >{\centering\arraybackslash}p{1.2cm}
}

\toprule[1.2pt]
\rowcolor{bleuprimaire}
\textcolor{white}{\bfseries Réf.} &
\textcolor{white}{\bfseries Compétence / Question} &
\textcolor{white}{\bfseries Corrigé détaillé \& critères d'évaluation} &
\textcolor{white}{\bfseries Barème} \\
\midrule[1pt]
\endfirsthead

\toprule[1.2pt]
\rowcolor{bleuprimaire}
\textcolor{white}{\bfseries Réf.} &
\textcolor{white}{\bfseries Compétence / Question} &
\textcolor{white}{\bfseries Corrigé détaillé \& critères d'évaluation} &
\textcolor{white}{\bfseries Barème} \\
\midrule[1pt]
\endhead

% ==================== EXERCICE 1 ====================
\rowcolor{bleusecondaire!15}
\multicolumn{4}{l}{
    \bfseries\textcolor{bleuprimaire}{\large Exercice 1 \hspace{0.2cm}\textbullet\hspace{0.2cm} Analyse et suites numériques} \hfill \bfseries 4 points
} \\* \midrule[0.6pt]

\textbf{1.1.a} &
Calculs de termes $u_1$ et $u_2$ &
\textbullet\ \subpt{0,25} $u_1 = f(0) = \frac{3(0)+1}{0+3} = \frac{1}{3}$.\newline
\textbullet\ \subpt{0,25} $u_2 = f\left(\frac{1}{3}\right) = \frac{3(1/3)+1}{1/3 + 3} = \frac{2}{10/3} = \frac{3}{5}$. &
\pts{0,50} \\
\cmidrule(lr){2-4}
\textbf{1.1.b} &
Preuve par récurrence &
\textbullet\ \subpt{0,25} Initialisation : $u_0 = 0$ donc $0 \le u_0 \le 3$ (vrai).\newline
\textbullet\ \subpt{0,50} Hérédité : Supposons $0 \le u_n \le 3$. $f$ est croissante sur $[0;3]$ donc $f(0) \le f(u_n) \le f(3)$. Or $f(0) = 1/3 \ge 0$ et $f(3) = 10/6 = 5/3 \le 3$. Donc $0 \le u_{n+1} \le 3$.\newline
\textbullet\ \subpt{0,25} Conclusion : La propriété est vraie pour tout $n \in \mathbb{N}$. &
\pts{1,00} \\
\cmidrule(lr){2-4}
\textbf{1.2.a} &
Variations de la fonction $f$ &
\textbullet\ \subpt{0,50} $f'(x) = \frac{3(x+3) - 1(3x+1)}{(x+3)^2} = \frac{3x+9-3x-1}{(x+3)^2} = \frac{8}{(x+3)^2}$. Pour tout $x \in [0;3]$, $f'(x) > 0$, donc $f$ est strictement croissante sur $[0;3]$. &
\pts{0,50} \\
\cmidrule(lr){2-4}
\textbf{1.2.b} &
Expression de $u_{n+1} - u_n$ &
\textbullet\ \subpt{0,50} $u_{n+1} - u_n = \frac{3u_n+1}{u_n+3} - u_n = \frac{3u_n+1 - u_n^2 - 3u_n}{u_n+3} = \frac{1 - u_n^2}{u_n+3} = \frac{(1-u_n)(1+u_n)}{u_n+3} = \frac{(3-u_n)(1+u_n)}{u_n+3}$. &
\pts{0,50} \\
\cmidrule(lr){2-4}
\textbf{1.2.c} &
Sens de variation de $(u_n)$ &
\textbullet\ \subpt{0,50} Puisque $0 \le u_n \le 3$, on a $3-u_n \ge 0$, $1+u_n > 0$ et $u_n+3 > 0$. Donc $u_{n+1} - u_n \ge 0$. La suite $(u_n)$ est donc croissante. &
\pts{0,50} \\
\cmidrule(lr){2-4}
\textbf{1.3} &
Convergence et limite &
\textbullet\ \subpt{0,25} $(u_n)$ est croissante et majorée par 3, donc elle converge vers $\ell$ par le théorème de convergence monotone.\newline
\textbullet\ \subpt{0,25} Par continuité de $f$, $\ell = \frac{3\ell+1}{\ell+3} \implies \ell(\ell+3) = 3\ell+1 \implies \ell^2 = 1$.\newline
\textbullet\ \subpt{0,25} Comme $u_n \ge 0$, $\ell \ge 0$, donc $\ell = 1$. &
\pts{0,75} \\

\midrule[1pt]

% ==================== EXERCICE 2 ====================
\rowcolor{bleusecondaire!15}
\multicolumn{4}{l}{
    \bfseries\textcolor{bleuprimaire}{\large Exercice 2 \hspace{0.2cm}\textbullet\hspace{0.2cm} Géométrie dans l'espace} \hfill \bfseries 5 points
} \\* \midrule[0.6pt]

\textbf{2.1.a} &
Non-alignement des points &
\textbullet\ \subpt{0,50} $\vect{AB}(2\,; -2\,; 1)$ et $\vect{AC}(0\,; 2\,; -1)$. Leurs coordonnées ne sont pas proportionnelles, donc $A, B, C$ ne sont pas alignés et définissent un plan. &
\pts{0,50} \\
\cmidrule(lr){2-4}
\textbf{2.1.b} &
Orthogonalité de $\vect{n}$ &
\textbullet\ \subpt{0,25} $\vect{n} \cdot \vect{AB} = 1(2) + 1(-2) + 2(1) = 0$.\newline
\textbullet\ \subpt{0,25} $\vect{n} \cdot \vect{AC} = 1(0) + 1(2) + 2(-1) = 0$. Donc $\vect{n}$ est normal au plan $(ABC)$. &
\pts{0,50} \\
\cmidrule(lr){2-4}
\textbf{2.1.c} &
Équation du plan $(ABC)$ &
\textbullet\ \subpt{0,50} $x + y + 2z + d = 0$. En passant par $A(1;2;1) : 1 + 2 + 2(1) + d = 0 \implies d = -5$. Équation : $x + y + 2z - 5 = 0$. &
\pts{0,50} \\
\cmidrule(lr){2-4}
\textbf{2.2.a} &
Direction de $(AD)$ &
\textbullet\ \subpt{0,50} $\vect{AD}(1\,; 0\,; 2)$. Le vecteur directeur $\vect{AD}$ n'est pas colinéaire au vecteur normal $\vect{n}(1\,; 1\,; 2)$, donc $(AD)$ n'est pas orthogonale au plan $(ABC)$. &
\pts{0,50} \\
\cmidrule(lr){2-4}
\textbf{2.2.b} &
Représentation paramétrique &
\textbullet\ \subpt{0,50} $(AD) : \begin{cases} x = 1 + t \\ y = 2 \\ z = 1 + 2t \end{cases}, t \in \mathbb{R}$. &
\pts{0,50} \\
\cmidrule(lr){2-4}
\textbf{2.2.c} &
Coordonnées du projeté $H$ &
\textbullet\ \subpt{0,50} En injectant la droite dans l'équation du plan : $(1+t) + 2 + 2(1+2t) - 5 = 0 \implies 5t = 0 \implies t = 0$. Donc $H = A(1\,; 2\,; 1)$. &
\pts{0,50} \\
\cmidrule(lr){2-4}
\textbf{2.3} &
Distance et volume &
\textbullet\ \subpt{0,25} Distance $D$ à $(ABC) = AH = \sqrt{(2-1)^2 + (2-2)^2 + (3-1)^2} = \sqrt{5}$.\newline
\textbullet\ \subpt{0,50} Aire de $ABC = \frac{1}{2}\|\vect{AB} \wedge \vect{AC}\| = \sqrt{5}$.\newline
\textbullet\ \subpt{0,25} Volume $= \frac{1}{3} \times \text{Aire} \times \text{hauteur} = \frac{5}{3}$. &
\pts{1,00} \\

\midrule[1pt]

% ==================== EXERCICE 3 ====================
\rowcolor{bleusecondaire!15}
\multicolumn{4}{l}{
    \bfseries\textcolor{bleuprimaire}{\large Exercice 3 \hspace{0.2cm}\textbullet\hspace{0.2cm} Probabilités} \hfill \bfseries 6 points
} \\* \midrule[0.6pt]

\textbf{3.1} &
Arbre pondéré &
\textbullet\ \subpt{0,50} Arbre à deux niveaux : $C$ (0,9) et $\bar{C}$ (0,1) ; sachant $C$, $T$ (0,95) et $\bar{T}$ (0,05) ; sachant $\bar{C}$, $T$ (0,02) et $\bar{T}$ (0,98). &
\pts{0,50} \\
\cmidrule(lr){2-4}
\textbf{3.2.a} &
Probabilité $P(C \cap T)$ &
\textbullet\ \subpt{0,50} $P(C \cap T) = P(C) \times P_C(T) = 0,9 \times 0,95 = 0,855$. &
\pts{0,50} \\
\cmidrule(lr){2-4}
\textbf{3.2.b} &
Probabilité totale de $T$ &
\textbullet\ \subpt{0,50} $P(T) = P(C \cap T) + P(\bar{C} \cap T) = 0,855 + 0,1 \times 0,02 = 0,875$. &
\pts{0,50} \\
\cmidrule(lr){2-4}
\textbf{3.3} &
Probabilité conditionnelle $P_T(\bar{C})$ &
\textbullet\ \subpt{0,50} $P_T(\bar{C}) = \frac{P(\bar{C} \cap T)}{P(T)} = \frac{0,002}{0,875} = \frac{2}{875} \approx 0,002\,29 \approx 0,002$. &
\pts{0,50} \\
\cmidrule(lr){2-4}
\textbf{3.4.a} &
Valeurs de $X$ &
\textbullet\ \subpt{0,50} Si accepté : $25 - 10 = 15$ \euro. Si refusé : $-10$ \euro. $X(\Omega) = \{-10\,; 15\}$. &
\pts{0,50} \\
\cmidrule(lr){2-4}
\textbf{3.4.b} &
Loi de probabilité de $X$ &
\textbullet\ \subpt{0,50} $P(X = 15) = 0,875$ ; $P(X = -10) = 0,125$. &
\pts{0,50} \\
\cmidrule(lr){2-4}
\textbf{3.4.c} &
Espérance mathématique &
\textbullet\ \subpt{0,50} $E(X) = 15 \times 0,875 + (-10) \times 0,125 = 11,875$ \euro. Bénéfice moyen de 11,875 \euro{} par composant. &
\pts{0,50} \\
\cmidrule(lr){2-4}
\textbf{3.5.a} &
Loi binomiale de $Y$ &
\textbullet\ \subpt{0,50} 50 épreuves indépendantes de Bernoulli de paramètre $p = 0,9$. $Y$ suit $\mathcal{B}(50\,; 0,9)$. &
\pts{0,50} \\
\cmidrule(lr){2-4}
\textbf{3.5.b} &
Calcul de $P(Y \ge 48)$ &
\textbullet\ \subpt{0,50} $P(Y \ge 48) = P(Y=48) + P(Y=49) + P(Y=50) \approx 0,032$. &
\pts{1,00} \\

\midrule[1pt]

% ==================== EXERCICE 4 ====================
\rowcolor{bleusecondaire!15}
\multicolumn{4}{l}{
    \bfseries\textcolor{bleuprimaire}{\large Exercice 4 \hspace{0.2cm}\textbullet\hspace{0.2cm} Fonctions, logarithme et convexité} \hfill \bfseries 5 points
} \\* \midrule[0.6pt]

\textbf{4.1.a} &
Limites de $g$ en $0$ et $+\infty$ &
\textbullet\ \subpt{0,50} $\lim_{x \to 0^+} g(x) = -\infty$ et $\lim_{x \to +\infty} g(x) = +\infty$. &
\pts{0,50} \\
\cmidrule(lr){2-4}
\textbf{4.1.b} &
Variations de $g$ &
\textbullet\ \subpt{0,50} $g'(x) = 2x + \frac{1}{x} > 0$ sur $]0 ; +\infty[$, donc $g$ est strictement croissante. &
\pts{0,50} \\
\cmidrule(lr){2-4}
\textbf{4.1.c} &
Bijection et encadrement $\alpha$ &
\textbullet\ \subpt{0,50} Par le corollaire du TVI, $g(x) = 0$ admet une unique solution $\alpha \approx 0,65$. &
\pts{0,50} \\
\cmidrule(lr){2-4}
\textbf{4.1.d} &
Signe de $g(x)$ &
\textbullet\ \subpt{0,50} $g(x) \le 0$ sur $]0 ; \alpha]$ et $g(x) \ge 0$ sur $[\alpha ; +\infty[$. &
\pts{0,50} \\
\cmidrule(lr){2-4}
\textbf{4.2.a} &
Limite de $f$ en $0$ &
\textbullet\ \subpt{0,50} $\lim_{x \to 0^+} f(x) = +\infty$. Asymptote verticale d'équation $x = 0$. &
\pts{0,50} \\
\cmidrule(lr){2-4}
\textbf{4.2.b} &
Limite en $+\infty$ et asymptote &
\textbullet\ \subpt{0,50} $\lim_{x \to +\infty} f(x) = +\infty$ et $\Delta : y = x$ est asymptote oblique en $+\infty$. &
\pts{0,50} \\
\cmidrule(lr){2-4}
\textbf{4.2.c} &
Expression de $f'(x)$ &
\textbullet\ \subpt{0,50} $f'(x) = \frac{g(x)}{x^2}$. &
\pts{0,50} \\
\cmidrule(lr){2-4}
\textbf{4.2.d} &
Variations de $f$ &
\textbullet\ \subpt{0,50} $f$ est décroissante sur $]0 ; \alpha]$ et croissante sur $[\alpha ; +\infty[$. &
\pts{0,50} \\
\cmidrule(lr){2-4}
\textbf{4.2.e} &
Point d'inflexion et tangente &
\textbullet\ \subpt{0,50} $f''(x) = \frac{3x - 2\ln(x)}{x^3}$, annulation en $x_0$ donnant le point d'inflexion $I$ et la tangente $\mathcal{T}$. &
\pts{0,50} \\

\midrule[1.2pt]
\rowcolor{bleuprimaire}
\multicolumn{3}{r}{\textcolor{white}{\bfseries\large TOTAL GÉNÉRAL}} &
\textcolor{white}{\bfseries\large 20 pts} \\
\bottomrule[1.2pt]

\end{longtable}

\end{document}